\documentclass[10pt,a4paper]{amsart}
\usepackage[margin=19mm]{geometry}
\usepackage{amsmath,amssymb,mathtools}
\usepackage[scr=rsfs,cal=euler]{mathalfa}
\usepackage{xfrac}
\usepackage[unicode,hidelinks,bookmarks=false]{hyperref}
\newcommand{\ie}{i.e.\ }
\newcommand{\cf}{cf.\ }
\newcommand{\resp}{respectively\ }
\newcommand{\Subsection}{\subsection}
\providecommand{\nobreakdash}{\nobreak\hbox{-}}
\setlength{\emergencystretch}{2em}
\multlinegap0pt
\usepackage{bm}
\usepackage{algorithm}\usepackage{algpseudocode}
\newtheorem{assumption}{Assumption}
\newtheorem{theorem}{Theorem}
\newcommand{\rd}{\mathrm{d}}
\title{Solving High-dimensional Nonlinear Space-time Fractional Diffusion PDEs via Deep-Picard Iteration}
\author{Zhijun Zeng, Zhitong Chen, Ling Qin, Yi Zhu}
\date{Source: arXiv:2605.00456v1 (CC BY 4.0)}
\begin{document}
\maketitle

\noindent\emph{Source and licence.} Solving High-dimensional Nonlinear Space-time Fractional Diffusion PDEs via Deep-Picard Iteration. Version arXiv:2605.00456v1 (CC BY 4.0), distributed by arXiv under the Creative Commons Attribution 4.0 International licence, http://creativecommons.org/licenses/by/4.0/, as recorded in the arXiv licence metadata for this exact version and shown on the abstract page https://arxiv.org/abs/2605.00456v1. Full licence notice: \texttt{licenses/arxiv-2605.00456-CC-BY-4.0.txt}.\par\smallskip

\noindent\textit{Editorial scope.} Counted unit: the article's theorem thm:weighted-contraction-final (source0.tex lines 497--510). The article's complete original proof of it is delivered with it (source0.tex lines 512--765, one \texttt{proof} environment of 254 lines). No mathematics is written, repaired or re-derived: the statement and the proof are the article's own lines, in the article's own notation, and no retained line is substituted or re-flowed.  Retained uncounted as the source-local context the proof needs: lines 335--341; lines 379--388; lines 407--416; lines 453--458; lines 463--487.  Nothing was omitted from any retained span, so the package carries no omission record.  No retained line contains a citation, so the wrapper carries no bibliography and no bibliography entry.  No retained line contains a figure environment, an image-inclusion command or drawing code, so no image is delivered and the package declares no asset dependency.  Editorial formatting only, in addition: an AMS \texttt{amsart} preamble that restates the article's own declarations and macros used by the retained lines, namely the environments \texttt{assumption}, \texttt{theorem} on separate counters, together with the standard packages those lines need.  The retained environments print in this document as Assumption 1, Theorem 1 in order of appearance, whereas the article prints the same items with the numbering of its own full document.  The complete native source of the article as distributed in the arXiv e-print for this exact version is bundled unchanged as \texttt{sources/199564/source0.tex}.
\begin{equation}\label{eq:stable-laplace}
\mathbb{E}\!\left[e^{-\lambda Y_s^\beta}\right]
=
e^{-s\lambda^\beta},
\qquad
\lambda>0,\quad s\ge0 .
\end{equation}

\begin{equation}\label{eq:terminal-payoff-prelim}
\mathcal{G}_{t,x}
=
u_0\!\bigl(X_{\tau_t}^x\bigr)\mathbf{1}_{\{\tau_t<\tau_\Omega(x)\}}
+
g\!\bigl(t_{\tau_\Omega(x)},X_{\tau_\Omega(x)}^x\bigr)
\mathbf{1}_{\{\tau_\Omega(x)\le\tau_t\}},
\qquad
t_{\tau_\Omega(x)}=t-Y_{\tau_\Omega(x)}^\beta .
\end{equation}

\begin{equation}\label{eq:T-operator-prelim}
(\mathcal{T}v)(t,x)
:=
\mathbb{E}_x\!\left[
\mathcal{G}_{t,x}
+
\int_0^{\tau(t,x)}
f\!\bigl(t_s,v(t_s,X_s^x),X_s^x\bigr)\,\rd s
\right],
\end{equation}

\begin{equation}\label{eq:weighted-norm-final}
\|v\|_\lambda
:=
\sup_{(t,x)\in[0,T]\times\overline{\Omega}}
e^{-\lambda t}|v(t,x)|.
\end{equation}

\begin{assumption}\label{ass:weighted-contraction}
The following conditions hold.
\begin{itemize}
    \item[(A1)] The nonlinearity $f:[0,T]\times\mathbb{R}\times\Omega\to\mathbb{R}$ is measurable and globally Lipschitz continuous in its second variable, uniformly in $(t,x)$, i.e.,
    \begin{equation}\label{eq:f-global-lip-final}
    |f(t,u,x)-f(t,v,x)|
    \le L_f |u-v|,
    \qquad
    \forall (t,x)\in[0,T]\times\Omega,\ \forall u,v\in\mathbb{R},
    \end{equation}
    for some constant $L_f>0$.
    
    \item[(A2)] The zero-level section of $f$ is bounded:
    \begin{equation}\label{eq:f-zero-bounded-final}
    M_f:=\sup_{(t,x)\in[0,T]\times\Omega}|f(t,0,x)|<\infty.
    \end{equation}
    
    \item[(A3)] The initial datum and boundary datum are bounded:
    \begin{equation}\label{eq:data-bounded-final}
    M_0:=\sup_{x\in\Omega}|u_0(x)|<\infty,
    \qquad
    M_g:=\sup_{(t,x)\in[0,T]\times\Omega^c}|g(t,x)|<\infty.
    \end{equation}
\end{itemize}
\end{assumption}

\begin{theorem}[Weighted contraction of the nonlinear Feynman--Kac operator]
\label{thm:weighted-contraction-final}
Suppose that Assumption~\ref{ass:weighted-contraction} holds. Then, for every $\lambda>0$, the operator $\mathcal{T}$ maps $\mathcal{X}_\lambda$ into itself. Moreover, for all $v,w\in\mathcal{X}_\lambda$,
\begin{equation}\label{eq:weighted-contraction-est-final}
\|\mathcal{T}v-\mathcal{T}w\|_\lambda
\le
\frac{L_f}{\lambda^\beta}\,\|v-w\|_\lambda .
\end{equation}
In particular, if
\begin{equation}\label{eq:lambda-condition-final}
\lambda^\beta>L_f,
\end{equation}
then $\mathcal{T}$ is a strict contraction on $\mathcal{X}_\lambda$.
\end{theorem}

\begin{proof}
The proof is divided into two steps.

\medskip
\noindent
\textbf{Step 1. Weighted Lipschitz estimate.}
Let $v,w\in\mathcal{X}_\lambda$ be arbitrary. Fix $(t,x)\in[0,T]\times\Omega$. Since the terminal payoff $\mathcal{G}_{t,x}$ is independent of the function argument of $\mathcal{T}$, we have
\begin{align}
(\mathcal{T}v)(t,x)-(\mathcal{T}w)(t,x)
&=
\mathbb{E}_x\!\left[
\int_0^{\tau(t,x)}
\Bigl(
f\!\bigl(t_s,v(t_s,X_s^x),X_s^x\bigr)
-
f\!\bigl(t_s,w(t_s,X_s^x),X_s^x\bigr)
\Bigr)\,\mathrm{d}s
\right].
\label{eq:weighted-proof-diff-final}
\end{align}
Taking absolute values and using \eqref{eq:f-global-lip-final}, we obtain
\begin{align}
|(\mathcal{T}v)(t,x)-(\mathcal{T}w)(t,x)|
&\le
L_f\,
\mathbb{E}_x\!\left[
\int_0^{\tau(t,x)}
|v(t_s,X_s^x)-w(t_s,X_s^x)|
\,\mathrm{d}s
\right].
\label{eq:weighted-proof-lip-final}
\end{align}
For every $0\le s<\tau(t,x)$, by the definition of the stopping time $\tau(t,x)=\tau_t(x)\wedge \tau_\Omega(x)$, we have
\[
t_s>0,
\qquad
X_s^x\in\Omega.
\]
Hence $(t_s,X_s^x)\in[0,T]\times\overline{\Omega}$, and by the definition of the weighted norm,
\begin{equation}\label{eq:weighted-pointwise-final}
|v(t_s,X_s^x)-w(t_s,X_s^x)|
\le
e^{\lambda t_s}\|v-w\|_\lambda .
\end{equation}
Substituting \eqref{eq:weighted-pointwise-final} into \eqref{eq:weighted-proof-lip-final}, and multiplying both sides by $e^{-\lambda t}$, yield
\begin{align}
e^{-\lambda t}|(\mathcal{T}v)(t,x)-(\mathcal{T}w)(t,x)|
&\le
L_f\|v-w\|_\lambda\,
\mathbb{E}_x\!\left[
\int_0^{\tau(t,x)}
e^{-\lambda t}e^{\lambda t_s}\,\mathrm{d}s
\right]
\nonumber\\
&=
L_f\|v-w\|_\lambda\,
\mathbb{E}_x\!\left[
\int_0^{\tau(t,x)}
e^{-\lambda(t-t_s)}\,\mathrm{d}s
\right].
\label{eq:weighted-proof-after-weight-final}
\end{align}
Since $t_s=t-Y_s^\beta$, we have
\(
t-t_s=Y_s^\beta,
\)
and therefore
\begin{equation}\label{eq:weighted-proof-before-tonelli-final}
e^{-\lambda t}|(\mathcal{T}v)(t,x)-(\mathcal{T}w)(t,x)|
\le
L_f\|v-w\|_\lambda\,
\mathbb{E}_x\!\left[
\int_0^{\tau(t,x)}
e^{-\lambda Y_s^\beta}\,\mathrm{d}s
\right].
\end{equation}
Since the integrand is nonnegative, Tonelli's theorem gives
\begin{align}
\mathbb{E}_x\!\left[
\int_0^{\tau(t,x)} e^{-\lambda Y_s^\beta}\,\mathrm{d}s
\right]
&=
\mathbb{E}_x\!\left[
\int_0^\infty \mathbf{1}_{\{s<\tau(t,x)\}} e^{-\lambda Y_s^\beta}\,\mathrm{d}s
\right]
\nonumber\\
&=
\int_0^\infty
\mathbb{E}_x\!\left[
\mathbf{1}_{\{s<\tau(t,x)\}} e^{-\lambda Y_s^\beta}
\right]
\,\mathrm{d}s
\nonumber\\
&\le
\int_0^\infty
\mathbb{E}\!\left[e^{-\lambda Y_s^\beta}\right]
\,\mathrm{d}s
\nonumber\\
&=
\int_0^\infty e^{-s\lambda^\beta}\,\mathrm{d}s
=
\frac{1}{\lambda^\beta},
\label{eq:weighted-proof-core-est-final}
\end{align}
where in the last step we used \eqref{eq:stable-laplace}. Combining \eqref{eq:weighted-proof-before-tonelli-final} and \eqref{eq:weighted-proof-core-est-final}, we obtain
\[
e^{-\lambda t}|(\mathcal{T}v)(t,x)-(\mathcal{T}w)(t,x)|
\le
\frac{L_f}{\lambda^\beta}\,\|v-w\|_\lambda.
\]
Taking the supremum over $(t,x)\in[0,T]\times\Omega$ gives
\[
\|\mathcal{T}v-\mathcal{T}w\|_\lambda
\le
\frac{L_f}{\lambda^\beta}\,\|v-w\|_\lambda,
\]
which proves \eqref{eq:weighted-contraction-est-final}.

\medskip
\noindent
\textbf{Step 2. The mapping property $\mathcal{T}(\mathcal{X}_\lambda)\subset\mathcal{X}_\lambda$.}
Let $v\in\mathcal{X}_\lambda$ and fix $(t,x)\in[0,T]\times\Omega$. By \eqref{eq:T-operator-prelim},
\begin{align}
|(\mathcal{T}v)(t,x)|
&\le
\mathbb{E}_x|\mathcal{G}_{t,x}|
+
\mathbb{E}_x\!\left[
\int_0^{\tau(t,x)}
\left|f\!\bigl(t_s,v(t_s,X_s^x),X_s^x\bigr)\right|
\,\mathrm{d}s
\right].
\label{eq:mapping-step1-final}
\end{align}
By \eqref{eq:f-global-lip-final} and \eqref{eq:f-zero-bounded-final},
\[
|f(t_s,v(t_s,X_s^x),X_s^x)|
\le
L_f|v(t_s,X_s^x)|+M_f.
\]
Hence
\begin{align}
|(\mathcal{T}v)(t,x)|
&\le
\mathbb{E}_x|\mathcal{G}_{t,x}|
+
L_f\,
\mathbb{E}_x\!\left[
\int_0^{\tau(t,x)}
|v(t_s,X_s^x)|\,\mathrm{d}s
\right]
+
M_f\,\mathbb{E}_x[\tau(t,x)].
\label{eq:mapping-step2-rough-final}
\end{align}
Instead of estimating the last term by $\mathbb{E}[\tau(t,x)]$, we exploit the same exponential weight as in Step~1. Multiplying \eqref{eq:mapping-step1-final} by $e^{-\lambda t}$, and using \eqref{eq:weighted-norm-final}, we obtain
\begin{align}
e^{-\lambda t}|(\mathcal{T}v)(t,x)|
&\le
e^{-\lambda t}\mathbb{E}_x|\mathcal{G}_{t,x}|
+
L_f\|v\|_\lambda
\mathbb{E}_x\!\left[
\int_0^{\tau(t,x)}
e^{-\lambda(t-t_s)}\,\mathrm{d}s
\right]
\nonumber\\
&\qquad
+
M_f\,
\mathbb{E}_x\!\left[
\int_0^{\tau(t,x)} e^{-\lambda t}\,\mathrm{d}s
\right].
\label{eq:mapping-step3-final}
\end{align}
We next estimate the three terms on the right-hand side.

First, by \eqref{eq:data-bounded-final} and the definition \eqref{eq:terminal-payoff-prelim},
\[
|\mathcal{G}_{t,x}|
\le
M_0\,\mathbf{1}_{\{\tau_t<\tau_\Omega(x)\}}
+
M_g\,\mathbf{1}_{\{\tau_\Omega(x)\le \tau_t\}}.
\]

Therefore,
\begin{equation}\label{eq:terminal-est-final}
e^{-\lambda t}\mathbb{E}_x|\mathcal{G}_{t,x}|
\le
\max\{M_0,M_g\}.
\end{equation}

Second, the same estimate as in \eqref{eq:weighted-proof-core-est-final} yields
\begin{equation}\label{eq:mapping-second-est-final}
L_f\|v\|_\lambda
\mathbb{E}_x\!\left[
\int_0^{\tau(t,x)}
e^{-\lambda(t-t_s)}\,\mathrm{d}s
\right]
\le
\frac{L_f}{\lambda^\beta}\|v\|_\lambda.
\end{equation}

Third, for $0\le s<\tau(t,x)$ we have $t_s\ge 0$, and hence
\[
e^{-\lambda t}
=
e^{-\lambda t_s}e^{-\lambda(t-t_s)}
\le
e^{-\lambda(t-t_s)}
=
e^{-\lambda Y_s^\beta}.
\]
Thus,
\begin{align}
M_f\,
\mathbb{E}_x\!\left[
\int_0^{\tau(t,x)} e^{-\lambda t}\,\mathrm{d}s
\right]
&\le
M_f\,
\mathbb{E}_x\!\left[
\int_0^{\tau(t,x)} e^{-\lambda Y_s^\beta}\,\mathrm{d}s
\right]
\le
\frac{M_f}{\lambda^\beta}.
\label{eq:mapping-third-est-final}
\end{align}
Combining \eqref{eq:mapping-step3-final}--\eqref{eq:mapping-third-est-final}, we arrive at
\[
e^{-\lambda t}|(\mathcal{T}v)(t,x)|
\le
\max\{M_0,M_g\}
+
\frac{L_f}{\lambda^\beta}\|v\|_\lambda
+
\frac{M_f}{\lambda^\beta}.
\]
Taking the supremum over $(t,x)\in[0,T]\times\Omega$, we obtain
\begin{equation}\label{eq:mapping-bound-final}
\|\mathcal{T}v\|_\lambda
\le
\max\{M_0,M_g\}
+
\frac{L_f}{\lambda^\beta}\|v\|_\lambda
+
\frac{M_f}{\lambda^\beta}
<\infty.
\end{equation}
Hence $\mathcal{T}v\in\mathcal{X}_\lambda$, i.e., $\mathcal{T}$ maps $\mathcal{X}_\lambda$ into itself.

The proof is complete.
\end{proof}

\end{document}
